% 图契约：进入、完全进入、离开时的磁通变化；图中的端口、几何、箭头和方程对应本题。
% 版面：沿用前三模块的几何图+推导分区；目标印宽17cm，简图至中等密度。
% 布局：左侧真实结构，右侧关键关系；真实电路边不作为架构流程箭头处理。
\documentclass[tikz,border=7pt]{standalone}
\usepackage{fontspec,xeCJK,amsmath,amssymb}
\IfFontExistsTF{Times New Roman}{\setmainfont{Times New Roman}}{\setmainfont{TeX Gyre Termes}}
\IfFontExistsTF{Heiti SC}{\setCJKmainfont[AutoFakeBold=1.5]{Heiti SC}}{\setCJKmainfont[AutoFakeBold=1.5]{FandolHei-Regular}}
\usetikzlibrary{arrows.meta,bending,calc,patterns,positioning,decorations.pathmorphing}
\definecolor{ink}{HTML}{243745}
\definecolor{blue}{HTML}{275DAB}
\definecolor{green}{HTML}{177C63}
\definecolor{red}{HTML}{BD4935}
\definecolor{orange}{HTML}{AA6C16}
\definecolor{purple}{HTML}{715393}
\tikzset{every node/.style={font=\fontsize{9.3}{12.5}\selectfont,text=ink},
arrow/.style={-{Stealth[length=5pt 1.5,width'=0pt 3]},shorten >=2pt,shorten <=1pt,line width=1pt},
arrow short/.style={-{Stealth[length=3pt 1.0,width'=0pt 2.5]},shorten >=1pt,shorten <=0pt,line width=.8pt},
axis/.style={arrow,draw=ink!60,line width=.6pt},
vec/.style={arrow,draw=blue},
wire/.style={draw=ink!75,line width=.8pt},
guide/.style={dashed,draw=ink!45,line width=.5pt},
card/.style={rounded corners=2pt,draw=blue!65,fill=blue!4,align=center,text width=6.5cm,inner xsep=9pt,inner ysep=7pt},
title/.style={font=\fontsize{10.5}{14}\selectfont\bfseries,text width=15cm,align=center},
label/.style={font=\fontsize{8.4}{11}\selectfont,align=center},
note/.style={label,text width=6.1cm}}
\begin{document}
\begin{tikzpicture}[x=1cm,y=1cm]
\path[use as bounding box] (-8,-3.5) rectangle(8,3.65);
\node[title] at (0,3.25) {M5-078a\quad 进入、完全进入、离开时的磁通变化};
\begin{scope}[shift={(-4.05,0)}]\fill[blue!4](-.6,1.15)rectangle(1.5,2.25);\draw[guide,blue](-.6,1.15)--(-.6,2.25)(1.5,1.15)--(1.5,2.25);\draw[wire,draw=ink](-1.4,2.1)--(-0.975,2.1) (-0.42499999999999993,2.1)--(0.0,2.1);\begin{scope}[shift={(-0.7,2.1)},rotate=0.0]\draw[draw=ink,fill=white,line width=.8pt](-0.275,-.085)rectangle(0.275,.085);\end{scope}
\draw[wire](-1.4,2.1)--(-1.4,1.2999999999999998)--(0.0,1.2999999999999998)--(0.0,2.1);\node[left]at(-1.85,1.7){进入};\draw[vec,green](2.35,1.7)--(2.95,1.7);\draw[arrow short,red](0.0,1.4)--(0.0,2.0);\fill[blue!4](-.6,-0.55)rectangle(1.5,0.55);\draw[guide,blue](-.6,-0.55)--(-.6,0.55)(1.5,-0.55)--(1.5,0.55);\draw[wire,draw=ink](-0.2,0.4)--(0.22499999999999998,0.4) (0.775,0.4)--(1.2,0.4);\begin{scope}[shift={(0.5,0.4)},rotate=0.0]\draw[draw=ink,fill=white,line width=.8pt](-0.275,-.085)rectangle(0.275,.085);\end{scope}
\draw[wire](-0.2,0.4)--(-0.2,-0.4)--(1.2,-0.4)--(1.2,0.4);\node[left]at(-1.85,0){全入};\draw[vec,green](2.35,0)--(2.95,0);\fill[blue!4](-.6,-2.25)rectangle(1.5,-1.15);\draw[guide,blue](-.6,-2.25)--(-.6,-1.15)(1.5,-2.25)--(1.5,-1.15);\draw[wire,draw=ink](0.8,-1.2999999999999998)--(1.225,-1.2999999999999998) (1.775,-1.2999999999999998)--(2.2,-1.2999999999999998);\begin{scope}[shift={(1.5,-1.2999999999999998)},rotate=0.0]\draw[draw=ink,fill=white,line width=.8pt](-0.275,-.085)rectangle(0.275,.085);\end{scope}
\draw[wire](0.8,-1.2999999999999998)--(0.8,-2.1)--(2.2,-2.1)--(2.2,-1.2999999999999998);\node[left]at(-1.85,-1.7){离开};\draw[vec,green](2.35,-1.7)--(2.95,-1.7);\draw[arrow short,red](0.8,-2.0)--(0.8,-1.4);\node[label]at(.5,2.55){固定磁区（浅蓝），$B$ 向纸内};\node[label]at(0,-2.65){绿色为整体速度；红箭头为电流方向};\end{scope}\node[card,draw=blue!65,fill=blue!4]at(3.9,1.8){进入：纸内磁通增加，电流逆时针};
\node[card,draw=green!65,fill=green!4]at(3.9,0){全入：磁通恒定，忽略自感时电流为零};
\node[card,draw=orange!65,fill=orange!4]at(3.9,-1.8){离开：纸内磁通减少，电流顺时针};

\end{tikzpicture}\end{document}